\subsection{UNIFORMITY AND COMPLETION OF A COMMUTATIVE TOPOLOGICAL GROUP}

We have already remarked that the left and right uniformities coincide on a commutative topological group G; whenever we speak of the uniformity of G, it is this unique uniformity to which we refer.

THEOREM 2. Let G be a commutative topological group. Then the functions \(x^{-1}\) and xy are uniformly continuous on G and G × G respectively. Moreover G admits a Hausdorff completion \(\hat{G}\) , and \(\hat{G}\) is commutative.

The uniform continuity of \(x^{-1}\) follows from Proposition 2 of no. 1, and that of xy from Proposition 3 of no. 1, since \((x, y) \to xy\) is a continuous homomorphism of \(G \times G\) into G. If G is Hausdorff, it satisfies the condition of Theorem 1 of no. 4 (as does every Hausdorff group whose left and right uniformities coincide); moreover the functions xy and yx are equal on \(\hat{G} \times \hat{G}\) by the principle of extension of identities; hence the second part of the theorem, by considering in the general case the Hausdorff group associated with G.

It follows in particular from this theorem that if \(f\) and \(g\) are two uniformly continuous mappings of a uniform space \(\mathbf{X}\) into a commutative group \(G\) , written additively, then the functions \(-f\) and \(f + g\) are uniformly continuous.

PROPOSITION 9. Let G be a commutative group, and let \(\mathfrak{T}_1\) , \(\mathfrak{T}_2\) be two Hausdorff topologies compatible with the group structure of G. Suppose that \(\mathfrak{T}_1\) is finer than \(\mathfrak{T}_2\) and that there is a fundamental system of neighbourhoods of o for \(\mathfrak{T}_1\) which are closed for \(\mathfrak{T}_2\) . Let \(\mathrm{G}_1, \mathrm{G}_2\) be the completions of G with respect to the topologies \(\mathfrak{T}_1, \mathfrak{T}_2\) respectively, and let \(f: \mathrm{G}_1 \to \mathrm{G}_2\) be the continuous homomorphism which extends the identity mapping of G (no. 3, Proposition 5). Then \(f\) is injective.

Suppose that G is written additively. Let \(\mathfrak{U}_{1}\) be the uniformity on G corresponding (no. 1) to the topology \(\mathfrak{T}_{1}\) : it will suffice to show that if F and \(F'\) are two minimal Cauchy filters (Chapter II, § 3, no. 2) with respect to \(\mathfrak{U}_{1}\) , which converge in \(G_{2}\) to the same point a, then \(F = F'\) (Chapter II, § 3, no. 7). For this, it is enough to show that \(F \cap F'\) is a Cauchy filter with respect to \(\mathfrak{U}_{1}\) . Let V be a neighbourhood of o in G with respect to \(\mathfrak{T}_{1}\) , such that V is closed in \(\mathfrak{T}_{2}\) , and let W be a symmetric neighbourhood of o in G with respect to \(\mathfrak{T}_{1}\) , such that \(W + W \subset V\) . By hypothesis, there is a \(W_{d}\) -small set M (resp. \(M'\) ) in F (resp. \(F'\) ); if \(x \in M\) and \(y \in M\) we have \(y - x \in W\) , i.e.~\(y \in x + W\) . If \(\overline{W}\) and \(\overline{V}\) are the closures of W and V in \(G_{2}\) , it follows that \(y \in x + \overline{W}\) , and therefore, since a is in the closure of M, that \(a \in x + \overline{W}\) for each \(x \in M\) . Similarly, \(a \in x' + \overline{W}\) for each \(x' \in M'\) , and hence \(x - x' \in \overline{W} + \overline{W}\) ; but since \((x, y) \to x + y\) is a continuous mapping of \(G_{2} \times G_{2}\) into \(G_{2}\) , we have \(\overline{W} + \overline{W} \subset \overline{W + W} \subset \overline{V}\) . It follows that if \(x \in M\) and \(x' \in M'\) , then \(x - x' \in \overline{V} \cap G = V\) , because V is closed in \(\mathfrak{T}_{2}\) ; and this completes the proof.

COROLLARY 1. Under the hypotheses of Proposition 9, if A is a subset of G which is a complete subspace with respect to the uniformity \(\mathfrak{U}_{2}\) corresponding to \(\mathfrak{T}_{2}\) , then A is also a complete subspace with respect to the uniformity \(\mathfrak{U}_{1}\) corresponding to \(\mathfrak{T}_{1}\) .

If \(A_{1}\) is the closure of A in \(G_{1}\) , then \(f(A_{1})\) is contained in the closure of A in \(G_{2}\) , which by hypothesis is equal to A. Since \(f(A)=A\) by definition and f is injective, we have \(A_{1}=A\) .

COROLLARY 2. Let G be a commutative group and let \(\mathfrak{T}_1, \mathfrak{T}_2\) be two topologies compatible with the group structure of G. Suppose that \(\mathfrak{T}_1\) is finer than \(\mathfrak{T}_2\) and that there is a fundamental system \(\mathfrak{B}\) of neighbourhoods of o with respect to \(\mathfrak{T}_1\) which are complete with respect to the uniformity \(\mathfrak{U}_2\) corresponding to \(\mathfrak{T}_2\) . Then G is complete with respect to the uniformity \(\mathfrak{U}_1\) corresponding to \(\mathfrak{T}_1\) .

The sets of \(\mathfrak{B}\) are closed in the topology \(\mathfrak{T}_{2}\) , hence complete for the uniformity \(\mathfrak{U}_{1}\) by Corollary 1; the result therefore follows from Proposition 4 of no. 3.

